% 第九卷 6.5.1.3 节：普通解码与推测解码的时间线（教学示意）。
% 目标模型一次前向记为宽 1.0 的块，草稿模型一步记为宽 0.25 的块；草稿长度 γ=4。
% 第一轮接受 3 个、拒绝第 4 个并补抽 1 个（产出 4 个）；第二轮全部接受并额外抽 1 个（产出 5 个）。
\begin{tikzpicture}[
  x=1cm,y=1cm,
  every node/.style={font=\figurefont\scriptsize,text=FigureInk},
  tgt/.style={rectangle,rounded corners=1.5pt,draw=FigureStroke,line width=0.6pt,fill=FigurePaper,minimum height=6mm,minimum width=9.4mm,inner sep=0pt},
  ver/.style={tgt,fill=FigureBlueFill,draw=FigureBlue},
  dft/.style={rectangle,rounded corners=1pt,draw=FigureYellow,line width=0.5pt,fill=FigureYellowFill,minimum height=6mm,minimum width=2.1mm,inner sep=0pt},
  acc/.style={rectangle,draw=white,line width=0.4pt,fill=FigureBlue,minimum size=1.8mm,inner sep=0pt},
  rej/.style={rectangle,draw=white,line width=0.4pt,fill=FigureRed,minimum size=1.8mm,inner sep=0pt},
  res/.style={rectangle,draw=white,line width=0.4pt,fill=FigureYellow,minimum size=1.8mm,inner sep=0pt},
  row/.style={font=\figurefont\bfseries\footnotesize,anchor=east},
  lab/.style={font=\figurefont\scriptsize,text=FigureMuted}
]
  % 普通解码。
  \node[row] at (-0.2,0) {普通解码};
  \foreach \k in {1,...,9} {
    \node[tgt] at ({\k-0.5},0) {目标};
    \node[acc] at ({\k-0.5},0.5) {};
  }
  \node[lab,anchor=west] at (9.1,0) {每次前向产出1个};

  % 推测解码。
  \node[row] at (-0.2,-1.5) {推测解码};
  \foreach \r in {0,1} {
    \foreach \i in {0,...,3} {\node[dft] at ({2*\r+0.125+\i*0.25},-1.5) {};}
    \node[ver] at ({2*\r+1.5},-1.5) {验证};
  }
  \node[lab,anchor=north] at (0.5,-1.85) {草稿$\times4$};
  \node[lab,anchor=north] at (2.5,-1.85) {草稿$\times4$};
  % 第一轮：接受3、拒绝1、补抽1。
  \foreach \i in {0,1,2} {\node[acc] at ({1.2+\i*0.2},-0.95) {};}
  \node[rej] at (1.8,-0.95) {};
  \node[res] at (1.8,-0.75) {};
  % 第二轮：全部接受，额外抽1。
  \foreach \i in {0,...,3} {\node[acc] at ({3.1+\i*0.2},-0.95) {};}
  \node[res] at (3.9,-0.95) {};
  \node[lab,anchor=west,align=left] at (4.45,-1.5) {两轮产出$4+5=9$个词元，\\耗时约为普通解码的4次前向};

  % 时间对齐线与图例。
  \draw[FigureMuted,dashed,line width=0.5pt] (4,0.8) -- (4,-2.1);
  \draw[-{Stealth[length=2mm]},FigureStroke,line width=0.6pt] (0,-2.35) -- (9.3,-2.35) node[lab,anchor=north east] {时间};
  \node[acc] at (0.2,-2.85) {}; \node[lab,anchor=west] at (0.35,-2.85) {接受或普通产出的词元};
  \node[rej] at (3.6,-2.85) {}; \node[lab,anchor=west] at (3.75,-2.85) {被拒绝的草稿词元};
  \node[res] at (6.4,-2.85) {}; \node[lab,anchor=west] at (6.55,-2.85) {目标模型补抽或额外抽取};
\end{tikzpicture}
